\begin{lemma}[Change lemma: ordered mixed location]
Fix a finite loopless multigraph with its exact edge set, a palette of
$k$ colors, and a distinguished edge $e_0$.  Let $(T,\varphi_n)$ be a
positive-rung extended Tashkinov tree.  Let $\mathcal B$ be a hierarchy
base and let
\[
 T_n=T_{n,0}\subset T_{n,1}\subset\cdots\subset
 T_{n,q}\subset T_{n,q+1}=T
\]
be a good hierarchy under $\varphi_n$, where $q$ is its last divider
index and
\[
 A=T_{n,q}^{*}=\noderef{HierarchyStarredStage}(H,q).
\]
Assume that $T$ has maximum property, that $V(T)$ is not elementary,
that $\varphi_n$ is proper and leaves exactly $e_0$ uncolored, and that
the rung number is positive.  These data form a
\noderef{HierarchyCounterexample} $C$.

Assume the six lower-rung assertions of the technical theorem, uniformly
for every positive-rung maximum ETT on the same graph, palette, and
distinguished edge, with arbitrary certified input and generating
colorings and with rung number smaller than that of $T$: lower-rung
elementarity, terminal interchangeability, the two PE stability and
stored-pair assertions, the terminal PE component assertion, the
defective-endpoint order assertion, and the mod-coloring assertion.
Assume also the same-rung induction hypothesis that every maximum ETT
with a good hierarchy having fewer than $q$ dividers is elementary.
Finally assume the large-color estimate
\[
 2n+11\leq |\overline\varphi_n(T_n)|.
\]

Put $\ell=|T|-|A|$, the number of exact tail additions above $A$, and let
\[
 T(y_{\ell-1})=T\big|_{\min\{|A|+(\ell-1),|T|\}}
\]
be the exact ordered prefix supplied by \noderef{TreeSequencePrefix}.
Assume (6.6): $V(T(y_{\ell-1}))$ is elementary under $\varphi_n$.
Assume moreover that $C$ is lexicographically minimal among all
\noderef{HierarchyCounterexample}s on the fixed graph, palette, and
$e_0$ which have exactly $q$ dividers, allowing their certified input
colorings, generating colorings, bases, ETTs, and good hierarchies to
vary.  Explicitly, no such counterexample $C'$ satisfies
\[
 \noderef{HierarchyCounterexampleLess}
  \bigl(\noderef{HierarchyCounterexampleMeasure}(C'),
        \noderef{HierarchyCounterexampleMeasure}(C)\bigr).
\]

Let $\alpha\ne\beta$ be colors missing on $T(y_{\ell-1})$.  Write
\[
 i_c=\noderef{HierarchyFirstMissingPosition}(\varphi_n,T,c),
 \qquad v_c=T.\mathit{vertex}(i_c).
\]
If at least one of $\alpha,\beta$ is not missing on $A$, assume both
\[
 v_\alpha\prec v_\beta                                      \tag{a}
\]
and the exact edge-identity exclusion
\[
 \alpha\notin
 \left\{c:\exists t\in\operatorname{Fin}|T|,
       |A|\leq t<i_\beta,
       \ \varphi_n(T.\mathit{edge}(t))=c\right\}.             \tag{b}
\]
Thus (b) is precisely
$\alpha\notin\varphi_n\langle T(v_\beta)-A\rangle$; it is not
an endpoint-only abbreviation.


Assume in addition that the first $\alpha$-missing vertex lies in the
last starred stage and the first $\beta$-missing vertex lies outside it:
\[
        v_\alpha\in V(A),\qquad v_\beta\notin V(A).
\]
Then the full two-color components through $v_\alpha$ and $v_\beta$
are equal in each of the following cases.
\begin{enumerate}
\item $q\geq1$, and either
$\alpha\in\overline\varphi_n(T_{n,q})$ or
$\{\alpha,\beta\}\cap D_{n,q}=\varnothing$.
\item $q=0$, and either
$\alpha\in\overline\varphi_n(T_n)$ or
$\{\alpha,\beta\}\cap D_n=\varnothing$.
\item $\alpha\in\overline\varphi_n(A)$ and the singleton
$\{\alpha\}$ is closed on $A$ under $\varphi_n$.
\end{enumerate}
In case 3 the colors are also $A$-interchangeable: at most one paper
$(\alpha,\beta)$-path component meets $A$, with an isolated vertex
component counted as a path.
\end{lemma}
\begin{proof}
Unpack $C$.  Its fields give positivity, the displayed
\noderef{TashkinovHierarchy}, its \noderef{GoodHierarchy} witness,
\noderef{MaximumProperty}, and non-elementarity.  Let
$K_x$ denote the full \noderef{TwoColorComponent} through the vertex
node $x$.  All components below therefore contain the actual edge nodes
of the multigraph; parallel edges are never identified by their endpoint
pairs.

We first discharge the zero-tail endpoint.  The nested-stage and terminal fields of
\noderef{TashkinovHierarchy}, together with the two cases of
\noderef{HierarchyStarredStage}, make $A$ an ordered initial segment of
$T$.  At index zero this is the `coreInsideFirst' field followed by the
stage chain, and at a positive index it is the corresponding ordinary
stage followed by that chain.  Hence
\[
                         |A|\leq |T|.                         \tag{0a}
\]

Suppose that $\ell=0$.  Since $\ell=|T|-|A|$ and (0a) gives
$|A|\leq |T|$, we have $|A|=|T|$.  Expanding
\noderef{TreeSequenceInitialSegment}, the already established inclusion
$A\subseteq T$ identifies every vertex entry and every exact edge entry;
because the lengths are equal, it identifies the ordered tree-sequences
themselves, so $A=T$.

For $\ell=0$, truncated subtraction gives $\ell-1=0$, and the length of the
displayed prefix in the statement is therefore
\[
 \min\{|A|+(\ell-1),|T|\}=\min\{|T|,|T|\}=|T|.
\]
The full-length instance of \noderef{TreeSequencePrefix} is entrywise
equal to $T$.  Thus (6.6), the assumed elementarity of this prefix, says
that $V(T)$ is elementary under $\varphi_n$, contradicting the
`notElementary' field of $C$ in
\noderef{HierarchyCounterexample}.  We have proved $\ell\ne0$, and hence
$\ell\geq1$.

It follows that $A$, every starred stage, and every ordinary hierarchy
stage through $T_{n,q}$ are ordered initial segments of
$T(y_{\ell-1})$.  Indeed, the prefix length is at least $|A|$ because
$\ell\geq1$, while \noderef{HierarchyBase},
\noderef{TashkinovHierarchy},
\noderef{HierarchyStarredStage}, and
\noderef{TreeSequenceInitialSegment} put all the named stages inside
$A$ entry by entry.  Restricting the pairwise disjoint missing-color
condition in \noderef{ColoredVertexSetElementary} therefore shows that
all these stages, including $A$, are elementary under $\varphi_n$.     \tag{1}

Because $\alpha$ and $\beta$ are missing on this elementary prefix,
the least-candidate clause of
\noderef{HierarchyFirstMissingPosition} makes $v_\alpha$ and
$v_\beta$ actual vertices of the prefix missing the indicated colors.
They are distinct unless the two first-missing positions name the same
vertex.  In that equality branch
$K_{v_\alpha}=K_{v_\beta}$ is immediate, but alternative 3 still requires
interchangeability on $A$.  If $q>0$, then
$A=T_{n,q}$ by \noderef{HierarchyStarredStage}, and the last-stage
instance of \noderef{step1} applies to the missing, closed color
$\alpha$.  If $q=0$ and the terminal kind is PE, conclusion (iii) of
\noderef{interchange} applies to the exact core $A$: reflexivity gives
both clauses of \noderef{CombinedStableColoring}, and the given
lower-rung package supplies its induction hypothesis.  Finally suppose
$q=0$ and the terminal kind is not PE.  Then
\noderef{HierarchyBase} identifies $A$ with $T_n$.  When the rung number
is greater than one, restrict the certified series through the preceding
rung; its final ladder closure is this exact $T_n$, maximum property
restricts, and the RE constructor of \noderef{TashkinovSeriesStep} gives
the required generating coloring while the SE constructor, together
with the reverse stability relation from \noderef{sc2}, gives the
stable-coloring premise.  Lower assertion (ii) therefore makes
$\alpha,\beta$ interchangeable on $T_n$.  At the one-rung endpoint the
same conclusion is the interchangeability clause of
\noderef{TechnicalTheoremBase}.  Thus the coincident-vertex branch proves
both every requested component equality and, under alternative 3, the
additional interchangeability conclusion.  We may henceforth assume
that the two vertices are distinct.  Elementarity says that $\beta$,
which is missing at $v_\beta$, is not missing at $v_\alpha$, and symmetrically
$\alpha$ is not missing at $v_\beta$.  By the complementary definitions
\noderef{MissingColors} and \noderef{PresentColors}, each vertex is
incident with an edge of the other color.  Properness makes it a
degree-one end in the two-color incidence graph.  Hence $K_{v_\alpha}$
and $K_{v_\beta}$ are nontrivial path components in the sense of
\noderef{TwoColorNontrivialPathComponent}.  This verifies the path-component
premise each time \noderef{TreeInterchangeable} is used below.           \tag{2}

We shall repeatedly use the following endpoint observation.  Under the
proper coloring $\varphi_n$, every two-color component has a simple
spanning path by \noderef{TwoColorComponentSpanningPath}; the incidence
degree is at most two by
\noderef{TwoColorIncidenceDegreeAtMostTwo}.  An end vertex of a
nontrivial such path is a graph vertex missing one of its two colors:
an edge-node has two incidence neighbors because the graph is loopless,
whereas at a graph vertex properness permits at most one edge of either
color, and degree one means that the other color is absent.  Thus a
nontrivial component which cannot cross the boundary of an elementary
set has both ends among the unique vertices of that set missing its two
colors.  Distinct components are vertex-disjoint by
\noderef{TwoColorComponentsDisjoint}.                                  \tag{3}

We next make explicit the preservation argument used whenever a
component through $v_\beta$ is changed.  Suppose $K_{v_\alpha}\ne
K_{v_\beta}$ and set
\[
 \sigma_n=\varphi_n/K_{v_\beta}.
\]
By \noderef{TwoColorComponentSwap} the swap changes only the exact
edge nodes of $K_{v_\beta}$ and only interchanges $\alpha$ and $\beta$.
It fixes the value $\mathsf{none}$, and
\noderef{TwoColorComponentSwapProper} makes $\sigma_n$ proper.
Consequently \noderef{IsOneEdgeUncolored} shows that precisely $e_0$
remains uncolored.  Since $v_\beta$ misses $\beta$ before the swap, it
misses $\alpha$ afterwards; the different component through
$v_\alpha$ is untouched and $v_\alpha$ still misses $\alpha$.
Therefore any certified prefix containing both vertices is
non-elementary under $\sigma_n$.                                       \tag{4}

Here is the exact smaller-counterexample test which will be cited in the
the location case below.  Let $B$ be a hierarchy stage or starred stage
which is an ordered prefix of the prefix $T(v_\beta)$, and suppose:
\begin{enumerate}
\item the swap is stable on $B$ with every control set required by its
hierarchy presentation, and every displayed edge identity of $B$ keeps
its value;
\item every displayed edge identity of $T(v_\beta)-B$ has the same
value under $\sigma_n$ and $\varphi_n$;
\item every vertex of $T(v_\beta)-v_\beta$ has the same missing-color
set under the two colorings;
\item the corresponding ETT, maximum property, hierarchy base, and all
five good-hierarchy fields have been verified under $\sigma_n$.
\end{enumerate}
Then these data and (4) form an actual
\noderef{HierarchyCounterexample} $C'$ with the same $q$ dividers.
The terminal stage is the exact prefix $T(v_\beta)$, not an abstract
tree having the same vertex set.

We also check its measure.  Put
$r=\noderef{HierarchyPathNumber}(H)$, keeping this path number distinct
from the full tail length $\ell$.  The vertex $v_\beta$ belongs to
$T(y_{\ell-1})$.  If it lies outside $A$, let $s$ be the length of the
tail of $T(v_\beta)$ after $A$.  Its position in the displayed prefix
gives $s\leq\ell-1<\ell$.  Let $r'$ be the path number of the truncated
hierarchy on $T(v_\beta)$.  Let $B=T_{n,q}$ be the ordinary last stage
used in the second coordinate of
\noderef{HierarchyCounterexampleMeasure}.  The old and truncated
hierarchies have the same $B$ and the same starred stage $A$.  Since
$B\preccurlyeq A\preccurlyeq T(v_\beta)\preccurlyeq T$, put
$d=|A|-|B|$; cancellation along these ordered initial segments gives
\[
 |T|-|B|=d+\ell,
 \qquad |T(v_\beta)|-|B|=d+s.                              \tag{4a}
\]

If $s<r$, the fallback-and-candidate definition in
\noderef{HierarchyPathNumber} gives $r'\leq s<r$.  If $r\leq s$, then
the original tail beginning at its $r$-th added vertex is a path, by the
defining candidate for $r$ (the fallback value $r=\ell$ is impossible
here because $r\leq s<\ell$); its
initial segment ending at $v_\beta$ is again a path, with a singleton
vertex allowed.  Hence $r$ is a candidate for the truncated hierarchy
and $r'\leq r$.  Thus either $r'<r$, or $r'=r$ and (4a), together with
$s<\ell$, shows that the second measure coordinate strictly decreases
from $d+\ell$ to $d+s$.  In both cases
\noderef{HierarchyCounterexampleLess} gives
\[
                  \mathfrak m(C')<_{\rm hier}\mathfrak m(C). \tag{5}
\]
The assumed minimality forbids $C'$.  Notice that (5) is applied only
after the actual base, hierarchy, maximum property, non-elementarity,
path number, and terminal-tail coordinate have all been checked.

\textbf{Location 3: $a\in V(A)$ and $b\notin V(A)$.}
The conditional assumptions again give (a)--(b).  The first-missing
definition and Tashkinov augmentation rule give
\[
 \alpha\notin\varphi_n\langle T(b)-A\rangle,
 \qquad
 \beta\notin\overline\varphi_n(T(b)-b),                       \tag{13}
\]
and no displayed edge of $T(b)-A$ has color $\beta$, except possibly
the exact first connecting edge when $q=0$, $A=T_n$, and
$\beta\in D_n$.  This is obtained step by step from
\noderef{TashkinovAugmentation}: an edge added before $b$ is colored by
a color missing on its preceding prefix, whereas $b$ is the first
vertex missing $\beta$.

Suppose first that $\alpha$ is closed on $A$.  If $q>0$, the
last-divider instance of \noderef{step1} makes $\alpha,\beta$
$A$-interchangeable.  If $q=0$, use conclusion (iii) of
\noderef{interchange} in PE, or the terminal-ladder conclusion of
\noderef{TechnicalTheoremBase} and the lower-rung induction interface in
the non-PE endpoint.  Thus $K_a$ is the only path component
meeting $A$.  If $K_a\ne K_b$, then $K_b$ is vertex-disjoint from $A$
by \noderef{TwoColorComponentsDisjoint}.  Swap $K_b$.  The outside-chain
lemma \noderef{OutsideKempeRecoloringStable}, (13), and the exact
construction (10)--(12) produce a same-divider counterexample on
$T(b)$, contradicting (5).  Therefore $K_a=K_b$.  Moreover the cited
interchangeability is exactly the additional conclusion required in
alternative 3.                                                        \tag{14}

We may henceforth assume
\[
        \alpha\text{ is not closed on }A,\qquad K_a\ne K_b,  \tag{15}
\]
and derive a contradiction under alternatives 1 and 2.

First let $q=0$.  Alternative 2 implies $\alpha\notin D_n$: this is
immediate from pair avoidance, while if $\alpha$ is missing on $T_n$ it
is removed from $D_n$ by \noderef{TashkinovStoredColors}.  If the last
extension were SE or RE, its constructor in
\noderef{TashkinovSeriesStep} would make every missing color outside
$D_n$ closed on $T_n=A$ (in SE use the selected stability and
\noderef{sc2}); this contradicts (15).  Hence the last extension is PE
and \noderef{HierarchyBase} gives the exact ordered union
\[
                         A=T_n\vee R_n.                       \tag{16}
\]

We now supply the elementarity premise required by the PE closure
argument.  If $n>1$, restrict the given ETT and maximum property through
rung $n-1$ and apply lower assertion (i).  If $n=1$, apply the
elementarity conclusion of \noderef{TechnicalTheoremBase}.  In either
case the pre-step coloring is elementary on $T_n$.  Thus
\noderef{ParallelCommonMissingClosed} applies with its genuine
pre-step-elementarity hypothesis.  If $a$ belonged to
$V(T_n)\cap V(R_n)$, then $\alpha$ would be common-missing on the two
constituents and the cited lemma would close it on their exact union,
contrary to (15).  Consequently
\[
 a\in V(T_n)\setminus V(R_n)\quad\hbox{or}\quad
 a\in V(R_n)\setminus V(T_n).                                \tag{17}
\]
This is where pre-step elementarity is used; PE alone has not been used
to infer closure.

In the first alternative of (17), conclusion (ii) of
\noderef{interchange} makes $K_a$ the only path component meeting
$T_n$; in the second, conclusion (i) does the same for $R_n$.  Thus the
different component $K_b$ is disjoint from the corresponding
constituent.  On the other constituent, (1), the fact that $a$ is its
unique possible $\alpha$-missing vertex in the core, and the fact that
$b$ is the first $\beta$-missing vertex outside the core show that every
vertex of $K_b$ in that constituent has both colors present.  A full
component swap therefore preserves its missing-color set.  Its incident
edge clause is preserved because any changed incident edge has color
$\alpha$ or $\beta$, neither of which is controlled there.  Hence the
swap is stable separately on $T_n$ with control $D_n$ and on $R_n$ with
empty control: it satisfies \noderef{CombinedStableColoring}.           \tag{18}

Equation (13) removes the possible first-edge exception: if it occurred,
the edge would have color $\beta\in D_n$, whereas alternative 2 and
$a\in A$ force either pair avoidance or $\alpha$ already missing on
$T_n$; in both cases the relevant conclusion of
\noderef{interchange} puts $K_b$ away from that exact connecting edge.
Apply \noderef{CombinedStableClosurePreserved} to (18).  It fixes every
displayed edge identity of both $T_n$ and $R_n$, and it makes the same
ordered tree $R_n$ a $\sigma_n$-closure of the same exact support prefix.
Thus every displayed identity of $T(b)$ retains its value.  Copy the
augmentation sequence edge by edge as in (11).  Apply \noderef{Td} to
the $T_n$ part of (18), obtain the corresponding ETT from
\noderef{ModColoring}, and transfer maximum property with
\noderef{MaximumPropertyTransport}.  The transported closure conclusion
and (16) supply the same base core.

For the zero-divider good hierarchy, pair size and disjointness are
literal and avoidance follows from the unchanged displayed prefix
edges.  We verify the two parts of the PE `parallelBase' field of
\noderef{GoodHierarchy} separately.  The vertexwise clauses of the two
constituent stable relations in (18) identify the missing-color sets on
$T_n$ and $R_n$ under $\sigma_n$ with their sets under $\varphi_n$.
They therefore identify the missing set of the exact union.  Together
with \noderef{HierarchyStoredColorsAt}, this keeps the stage-zero stored
index set and hence the copied union $\Gamma^0$ unchanged.  It also
preserves the cardinality
\[
 \left|\bigl(\overline\varphi_n(T_n)\cap
       \overline\varphi_n(R_n)\bigr)-\Gamma^0\right|\geq4
\]
with the identical inequality under $\sigma_n$.

It remains to transport non-minus-closedness; combined stability does
not assert equality on every boundary edge.  Unfold the original
`parallelBase' witness using \noderef{TreeColorSetMinusClosed},
\noderef{TreeColorSetClosed}, \noderef{BoundaryColorEdges}, and
\noderef{BoundaryEdges}.  There are a color
\[
 c\in\overline\varphi_n(A)-\Gamma^0
\]
and an exact boundary-edge identity $f$ of $A$ with
$\varphi_n(f)=c$.  If the component swap fixes $f$, the same $c$ and $f$
witness that $A$ is not $(\Gamma^0)^-$-closed under $\sigma_n$.

Suppose instead that $f$ is changed.  By
\noderef{TwoColorComponentSwap}, its edge-node lies in $K_b$ and
$c\in\{\alpha,\beta\}$.  Since $b$ is the first vertex missing $\beta$
and $b\notin A$, the color $\beta$ is not missing anywhere on $A$.
Thus $c\ne\beta$, and consequently
\[
             c=\alpha,\qquad \alpha\notin\Gamma^0,
             \qquad f\in K_b.                                \tag{19a}
\]
Assume for contradiction that the transported core $A$ is
$(\Gamma^0)^-$-closed under $\sigma_n$.  The missing-set equality from
(18), (19a), and the definitions just cited say that $\alpha$ is missing
and singleton-closed on $A$ under $\sigma_n$.  Apply conclusion (iii) of
\noderef{interchange}: properness and exact one-uncoloredness were proved
in (4), its combined-stability premise is exactly (18), and its
lower-rung premise is the stated six-assertion package.  Hence
$\alpha,\beta$ are $A$-interchangeable under $\sigma_n$.

On the other hand, swapping a full component changes every one of its
$\alpha$- and $\beta$-colored edge identities to the other color and
changes no membership in the two-color incidence graph.  By
\noderef{TwoColorComponentSwap} and
\noderef{TwoColorComponent}, the old supports $K_a$ and $K_b$ are
therefore the same two distinct full components under $\sigma_n$.
The definition \noderef{TwoColorNontrivialPathComponent} and (2) show that they
remain nontrivial path components.  The component $K_a$
meets $A$ at $a$.  The boundary identity $f$ in (19a) has exactly one
endpoint in $A$ by \noderef{BoundaryEdges}; incidence with that
inside endpoint shows that $K_b$ also meets $A$.  Their distinctness is
preserved, equivalently by \noderef{TwoColorComponentsDisjoint}.  These
two components contradict the just-established
\noderef{TreeInterchangeable}.  Therefore $A$ is not
$(\Gamma^0)^-$-closed under $\sigma_n$ after all.

This proves the entire PE base field.  There are no refresh or
positive-stage minus-closed fields when $q=0$.  Thus the actual candidate
satisfies every good-hierarchy field, is non-elementary by (4), and has
the measure comparison (5).  This contradicts minimality and settles
$q=0$.                                                                 \tag{19}

It remains to treat $q>0$.  Since $a\in A=T_{n,q}$, the color $\alpha$
is missing on $T_{n,q}$, so \noderef{rutcor} applies.  Its full-tail
prefix premise is exactly the ambient (6.6): both use
$A=T_{n,q}^{*}$, $\ell=|T|-|A|$, and the prefix of length
$\min\{|A|+(\ell-1),|T|\}$.  Its pre-step
elementarity premise is supplied exactly as in the paragraph preceding
(17), from lower assertion (i) or
\noderef{TechnicalTheoremBase}; hence this use also respects the PE
common-missing restriction.

Suppose first that some starred stage is $\alpha$-closed, and choose the
largest such index $r$.  By (15), $r<q$.  The first conclusion of
\noderef{rutcor} gives the exact edge-index exclusion
\[
 \alpha\notin
       \varphi_n\langle T_{n,q}-T_{n,r}^{*}\rangle.            \tag{20}
\]

Before using any interchangeability result, we prove the missing-color
premise which closedness alone does not supply.  We claim, by descending
induction, that $\alpha$ is missing on every starred stage
$T_{n,s}^{*}$ with $r\leq s\leq q$.  At $s=q$ the witness is
$a\in A=T_{n,q}^{*}$.  Now let $r<s\leq q$, assume $\alpha$ is missing
on $T_{n,s}^{*}=T_{n,s}$, and suppose it is not missing on
$T_{n,s-1}^{*}$.  A vertex witnessing missingness on $T_{n,s}$ then lies
in the exact difference
\[
 V(T_{n,s})-V(T_{n,s-1}^{*}).
\]
The stage is elementary by (1), so
\noderef{GoodHierarchyNewMissingClosed} makes $\alpha$ closed on
$T_{n,s}=T_{n,s}^{*}$.  This contradicts the maximal choice of $r$,
because $s>r$.  Thus missingness descends from $s$ to $s-1$; the last
step $s=r+1$ proves in particular
\[
          \alpha\in\overline\varphi_n(T_{n,r}^{*}).           \tag{20a}
\]

Choose a vertex $x$ of $T_{n,r}^{*}$ witnessing (20a).  That starred
stage is an ordered initial segment of $T$.  The least-candidate clause
of \noderef{HierarchyFirstMissingPosition} says that the displayed
position of $a=v_\alpha$ is no later than the displayed position of
$x$.  The entrywise initial-segment relation then puts $a$ itself in
$T_{n,r}^{*}$.                                                \tag{20b}

We may now invoke interchangeability with all of its premises present.
If $r>0$, (20a), the defining closedness of $r$, (1), the given lower
package, the fewer-divider induction, maximum property, and the large
estimate are exactly the hypotheses of \noderef{step1}.  If $r=0$ and
the terminal step is PE, conclusion (iii) of \noderef{interchange}
applies to the missing, singleton-closed color on the exact core, with
reflexive \noderef{CombinedStableColoring}.  If $r=0$ and the terminal
step is not PE, use the lower-rung terminal-ladder conclusion, or
\noderef{TechnicalTheoremBase} at the one-rung endpoint.  Hence the two
colors are $T_{n,r}^{*}$-interchangeable in every case.

By (20b), $K_a$ meets this stage.  Both $K_a$ and $K_b$ are nontrivial
path components by (2), and they are distinct by (15).  If $K_b$ also
met the stage, \noderef{TreeInterchangeable} would identify them.
Therefore $K_b$ is vertex-disjoint from $T_{n,r}^{*}$.

Apply \noderef{OutsideKempeRecoloringStable} now, before any call to
\noderef{Td}.  The terminal ladder tree $T_n$ is an ordered initial
segment of $T_{n,r}^{*}$ by \noderef{HierarchyBase} and
\noderef{TashkinovHierarchy}; hence the swap is outside both trees.  In
particular, $\sigma_n$ is proper, leaves exactly $e_0$ uncolored, is
\[
       (T_n,D_n,\varphi_n)\text{-stable},                     \tag{20c}
\]
and fixes every exact edge identity incident with $T_n$.
If the terminal kind is PE, its exact core $T_n\vee R_n$ is also an
ordered initial segment of $T_{n,r}^{*}$.  Applying the same outside
lemma to each constituent gives \noderef{CombinedStableColoring}, and
\noderef{CombinedStableClosurePreserved} transports the exact old
closure $R_n$ and the ordered union.

If $r>0$, truncate the current terminal augmentation sequence at the
exact tree $T_{n,r}^{*}=T_{n,r}$.  This is a same-rung ETT whose terminal
tree is $T_{n,r}$; it carries the first $r-1$ dividers, because its
terminal stage has index $(r-1)+1$.  Its ladder series, stored sets,
connecting-edge identities, and extension kinds are unchanged, so every
field of maximum property restricts.  Pair size, avoidance,
disjointness, unique refresh, the PE base field, and the minus-closed
fields through index $r-1$ are literal restrictions of the old good
hierarchy.  Since $K_b$ is disjoint from $T_{n,r}^{*}$, its swap is an
outside recoloring of this truncated ETT.  Apply
\noderef{HierarchyLEMStable} with the stated six lower-rung assertions.
It supplies the corresponding certified retained prefix, maximum
property, and the transported good-hierarchy fields through $r$.

If $r=0$, do not assert that the starred core is itself a same-rung ETT.
In the PE case $K_b$ is outside the exact core
$T_{n,0}^{*}=T_n\vee R_n$, so use the combined stability and transported
closure just recorded.
In the non-PE case the core is $T_n$, and the same outside-recoloring
lemma gives terminal-ladder stability directly.  Thus, before rebuilding
anything above the retained stage, we have (20c) and its incident-edge
equality in every case; for $r>0$ we additionally have the actual
corresponding retained ETT from \noderef{HierarchyLEMStable}, while for
$r=0$ we have the exact core and closure data needed to rebuild the same
base.                                                                   \tag{21}

Combine (20) with (13).  Their adjacent exact index intervals show that
$\alpha$ is unused on $T(b)-T_{n,r}^{*}$.  For $\beta$, inspect the
same exact identities rather than make a blanket assertion.  Every
ordinary edge after $T_{n,r}^{*}$ and through $b$ is introduced by
\noderef{TashkinovAugmentation}; its color is missing on the preceding
prefix.  Since $b$ is the first vertex missing $\beta$, none of these
ordinary identities has color $\beta$.  The only possible exception is
the manuscript's exact one: $r=0$, $T_{n,0}^{*}=T_n$, and the first
displayed identity is the SE/RE connecting seed $f_n$ with
$\beta\in D_n$.

In that exceptional configuration, unpack the terminal SE or RE
constructor of \noderef{TashkinovSeriesStep}.  Its exact
\noderef{TreeSequenceAdjoinEdge} first forms $T_n+f_n$, whose new entry
at index $|T_n|$ is the individual edge identity $f_n$; the subsequent
\noderef{TashkinovClosure} preserves that entry.  The strict first-stage
and common-terminal initial-segment fields of
\noderef{TashkinovHierarchy} therefore identify the first displayed
edge after $T_n$ with this same $f_n$.  In SE the constructor's
\noderef{MaximumDefectiveChoice} makes $f_n$ a boundary identity incident
with its chosen vertex of $T_n$; in RE the recorded revisiting
certificate in \noderef{RevisitingExtensionCertificate} makes $f_n$ a
boundary identity with its inside endpoint in $T_n$.  Thus $f_n$ is
incident with $T_n$.  If the exception is realized, its old edge
equation is $\varphi_n(f_n)=\beta$, and $\beta\in D_n$.  The edge clause
of \noderef{StableColoring} in the already established relation (20c),
equivalently the stronger incident-edge conclusion of
\noderef{OutsideKempeRecoloringStable}, gives
\[
                         \sigma_n(f_n)=\varphi_n(f_n).         \tag{20d}
\]
This conclusion precedes reconstruction; it has not been read from a
corresponding ETT whose construction would already require the seed.
Consequently (13), (20), the ordinary augmentation argument, and (20d)
show that every displayed edge identity in the appropriate PE seed or
exact SE/RE one-edge seed and the subsequent sequence to $T(b)$ has its
old value under $\sigma_n$.

The missing-color witnesses in that sequence are preserved as well.
Before $b$, no vertex misses $\beta$.  By elementarity (1), $a$ is the
only prefix vertex missing $\alpha$; the distinct component $K_b$ does
not contain $a$, while at each of its other vertices before $b$ both
colors are present and the swap merely exchanges them.  Vertices outside
$K_b$ are fixed.  Hence every old augmentation witness is still missing
the same edge color under $\sigma_n$, and the same exact augmentation
sequence is valid.

Only now apply \noderef{Td} to (20c), using the stated six lower-rung
assertions and the original maximum property.  It gives the required
\noderef{ModColoring}.  We also form the old-coloring presentation on
the exact prefix $T(b)$: truncate the original terminal
\noderef{TashkinovClosure} at its displayed length, retain the certified
series and maximum property, retain the same hierarchy base and stages
through $T_{n,q}$, and replace only the terminal hierarchy stage by
$T(b)$.  The last containment is strict because $b\notin A$, and every
entrywise and divider field is inherited from
\noderef{TashkinovHierarchy}.  Thus this is a positive-rung maximum ETT
on $T(b)$ under $\varphi_n$ with the same base and hierarchy.

If the terminal kind is PE, apply \noderef{HierarchySplitter} to that
exact old-coloring prefix presentation.  The combined stability
established after (20c), the exact transported closure, and the
just-verified Tashkinov augmentation route are its remaining hypotheses.
The lemma supplies an actual
\noderef{CorrespondingExtendedTashkinovTree} on $T(b)$, maximum property,
and a hierarchy with the same $R_n$, starred core, stages, and divider
positions.  If the terminal kind is SE or RE, the verified exact
one-edge seed and subsequent route make $T(b)$ eligible for the universal
clause of \noderef{ModColoring}; its corresponding certificate preserves
the non-PE kind and terminal ladder tree, so \noderef{HierarchyBase}
rebuilds with the common core $T_n$ and the hierarchy fields copy entry
by entry.  In this branch \noderef{MaximumPropertyTransport} supplies
maximum property.

Combine the resulting full certificate and hierarchy with the retained
good-hierarchy data in (21): for $r>0$ use the transported Gamma pairs
and fields supplied by \noderef{HierarchyLEMStable}, while for $r=0$ use
the transported PE closure or the common non-PE core $T_n$.  Thus the
same base, retained stages, Gamma pairs, and retained minus-closed fields
are available before the later stages are checked.

We check the good hierarchy after $r$, where (21) alone does not suffice.
The Gamma pairs themselves, their sizes, and disjointness are unchanged.
The stage missing sets are unchanged because $b$ is later than every
divider and every earlier vertex of $K_b$ has both colors present.
Avoidance and unique refresh are unchanged because the displayed
segment edges and stage-difference missing sets are unchanged.  For a
stage $j>r$, maximality of $r$ says that $\alpha$ is not closed there.
The descending induction (20a) already proves that it is missing there;
no closedness-to-missingness implication is being used.  The minus-closed
field of
\noderef{GoodHierarchy} would close any missing color outside its
reserved union.  Since $\alpha$ is not closed, it follows that
$\alpha$ lies in
$\bigcup_{\eta\in D_{n,j}}\Gamma^{j-1}_\eta$; $\beta$ is not missing on
that stage because its first missing vertex is $b$.  Hence changing an
$\alpha$- or $\beta$-colored boundary identity cannot violate the
minus-closed assertion: $\alpha$ is among the excluded reserved colors
and $\beta$ is not among the colors to which minus-closedness applies.
All other boundary colors are fixed.  This verifies the fifth field
stage by stage with the same Gamma pairs.  The parallel-base field lies
at or before $r$ and was transported in (21).                              \tag{22}

The corresponding ETT and maximum property just obtained, the data
(21)--(22), (4), and the copied terminal sequence constitute a
same-divider hierarchy counterexample on $T(b)$.  Its actual measure is
smaller by (5), contradicting minimality.

Finally suppose no starred stage is $\alpha$-closed.  Use every part of
the second conclusion of \noderef{rutcor}.  For every $1\leq j\leq q$,
\[
 \alpha\in\bigcup_{\eta\in D_{n,j}}\Gamma^{j-1}_\eta
       \subseteq\Gamma^{j-1};                                \tag{23}
\]
the last extension is PE; $a\in V(T_n)\setminus V(R_n)$; and
\[
 \alpha\notin
       \varphi_n\langle T_{n,q}-T_n\rangle.                   \tag{24}
\]
By (1), no vertex of $R_n$ misses $\alpha$, since such a vertex would be
different from $a$ in the elementary core.  No vertex of $R_n$ misses
$\beta$, because $b$ is its first missing vertex and lies outside $A$.
Conclusion (ii) of \noderef{interchange} makes $K_b$ disjoint from
$T_n$.  Consequently the component swap is stable on $T_n$ and,
separately, on $R_n$: on $T_n$ it changes no incident identity, and on
$R_n$ every encountered component vertex has both colors present, so
the swap preserves its missing set.  This gives the exact
\noderef{CombinedStableColoring} premise, not an unsupported assertion
that PE alone supplies closure.                                        \tag{25}

Equations (24), (13), and the PE-base edge interval cover
$T(b)-T_n$.  The closure sequence of $R_n$ and elementarity of the core
exclude $\alpha$ from each exact edge of $R_n-T_n$; otherwise the
augmentation supplying that edge would name a second core vertex missing
$\alpha$.  The same first-missing argument excludes $\beta$.  Apply
\noderef{CombinedStableClosurePreserved} to (25).  It fixes every
displayed identity of $T_n$ and $R_n$ and transports the exact closure.
Together with the stated exclusions, all displayed edge identities of
$T(b)$ are fixed.  Apply \noderef{Td} to
the first stability relation in (25), copy the exact PE closure and
terminal augmentation sequences, obtain the corresponding ETT from
\noderef{ModColoring}, and transport maximum property by
\noderef{MaximumPropertyTransport}.                                   \tag{26}

Copy the hierarchy and Gamma pairs.  Pair size, disjointness, avoidance,
unique refresh, and the PE base field are preserved by (25)--(26).
For every positive stage, (23) places $\alpha$ in the exact reserved
union excluded from its minus-closed requirement, while $\beta$ is not
missing before $b$; all other boundary colors are unchanged.  Hence the
same-Gamma minus-closed field holds stage by stage.  We have again
constructed the actual base, hierarchy, good witness, corresponding
maximum ETT, and non-elementarity witness required by
\noderef{HierarchyCounterexample}.  The smaller measure (5) contradicts
minimality.                                                            \tag{27}

All possibilities under (15) contradict minimality, so $K_a=K_b$ in
this mixed location under alternatives 1 and 2.  Under alternative 3,
$\alpha$ is closed on $A$, so (14) gives both the same component equality
and \noderef{TreeInterchangeable} on $A$.  By its definition, at most one
paper path component meets $A$, including an isolated vertex component.
These are exactly the asserted conclusions.
\end{proof}
